\section{Introdução}
\label{sec:intro}

Os grandes modelos de linguagem tornaram-se infraestrutura de trabalho da
administração pública em poucos anos. Um levantamento da OCDE sobre funções
centrais de governo documenta a adoção de IA generativa em formulação de
políticas e prestação de serviços, incluindo assistentes que ajudam assessores
do serviço público a recuperar informação e as fontes por trás
dela~\citep{oecd2024governing}. A política de poluição do ar é um domínio em que
isso tem consequência direta para a saúde pública. O material particulado fino
está entre os principais fatores de risco ambiental para mortalidade prematura
no mundo~\citep{gbd2019risk,who2021aqg}, e a carga recai de forma
desproporcional sobre o Sul Global. Quando um gestor ambiental pergunta a um
modelo sobre um padrão de qualidade do ar, um prazo de atendimento ou a
evidência por trás de uma intervenção, uma resposta imprecisa pode se propagar
para um parecer regulatório, uma decisão de compra ou uma resposta de emergência
durante um episódio de poluição.

Este artigo pergunta se os modelos respondem a essas questões com a mesma
confiabilidade para países do Sul Global e do Norte Global, e se a forma de
enquadrar a consulta muda a resposta. As duas premissas têm apoio. O viés
factual geográfico é mensurável: as taxas de erro são $1,5\times$ maiores para
países da África Subsaariana do que para a América do Norte em 20
modelos~\citep{moayeri2024worldbench}, as avaliações dos modelos sobre temas
subjetivos acompanham condições socioeconômicas até
$\rho=0,70$~\citep{manvi2024llm}, e os modelos privilegiam sistematicamente
afirmações sobre o Norte Global~\citep{mirza2024global}. A assimetria linguística
por trás disso é estrutural: os idiomas da maioria dos países do Sul Global
ocupam as classes mal servidas da hierarquia de recursos~\citep{joshi2020state},
e grandes corpora derivados da Internet ``super-representam pontos de vista
hegemônicos e codificam vieses potencialmente danosos a populações
marginalizadas''~\citep{bender2021stochastic}. A produção decolonial lê essa
assimetria como uma infraestrutura epistêmica desigual reproduzida pelos
sistemas de IA~\citep{mohamed2020decolonial}.

Três lacunas organizam o estudo. Nenhum benchmark visa informação regulatória
com gabarito oficial: as avaliações existentes usam recordação de indicadores do
Banco Mundial, avaliações subjetivas, checagem jornalística ou conhecimento
cultural cotidiano, e nenhuma mede as tarefas que um gestor delega, como
enunciar um padrão vinculante, recuperar uma concentração medida oficialmente,
resumir a evidência de carga de doença, identificar instrumentos legais e órgãos
executores ou recomendar uma resposta operacional. O enquadramento por persona
não foi examinado como alavanca experimental, embora os praticantes prefixem um
papel rotineiramente e os fornecedores o recomendem (Seção~\ref{sec:related});
se ele reduz a lacuna, ao direcionar para respostas conformes à regulação, ou a
amplia, ao convidar à fabricação confiante, não foi testado em desenho
pré-especificado. E o mecanismo de representação no corpus costuma ser afirmado
em vez de medido contra confundidores, e raramente é separado em seus dois
canais, o tamanho do corpus do idioma~\citep{xue2021mt5,abadji2022oscar} e a
amplitude da cobertura do país, distinção que importa aqui porque a lacuna
geográfica é medida em inglês.

\subsection*{Objetivos e escopo}

O estudo tem três objetivos delimitados. O primeiro é estabelecer um protocolo
de auditoria de informação regulatória. Seu objeto é a concordância entre a
resposta de um modelo e o valor publicado no registro oficial, para tarefas que
um órgão ambiental delega; ele não mede capacidade geral do modelo, qualidade de
raciocínio ou utilidade para outros fins. O protocolo é a contribuição que se
transfere, e o domínio é onde o demonstramos. O segundo é testar as mitigações
disponíveis a uma administração. Um órgão público diante de respostas ruins
sobre o próprio país pode declarar o papel oficial no prompt, perguntar no
idioma nacional ou rodar um modelo aberto com ajuste regional, e testamos as
três como fatores de desenho. O modelo regional que pudemos rodar é um ajuste
comunitário de uma base de 7B em português brasileiro, do tipo que um órgão
consegue baixar; o achado diz respeito a essa classe, e o protocolo liberado
permite testar qualquer outro candidato. Não testamos recuperação aumentada,
ajuste fino ou verificação com humano no circuito, cada um dos quais exige
capacidade de engenharia que as correções no nível do prompt não exigem. Todas as consultas
são respondidas a partir da memória paramétrica do modelo, sem busca na web nem
recuperação de documentos, e por isso o título fala em informação recordada, e
não recuperada. O terceiro objetivo é localizar onde a confiabilidade difere e,
até onde o desenho permite, por quê; testamos a explicação por representação no
corpus e não reivindicamos identificação causal.

Ao longo do texto, a unidade de análise é o modelo tal como servido, isto é,
pesos, fornecedor e infraestrutura de serviço alcançados por um endpoint fixado,
porque o gestor público interage com essa configuração e qualquer lacuna medida é
propriedade dela. Hipóteses, amostragem e plano de análise foram fixados antes
da coleta~\citep{nosek2018preregistration} para 15 países e depois estendidos a
25; o plano nunca foi depositado publicamente, portanto não reivindicamos
pré-registro, e o estudo é relatado como exploratório do início ao fim.

\subsection*{Contribuições deste trabalho}

\begin{enumerate}\setlength{\itemsep}{2pt}
\item \emph{Um protocolo de auditoria para informação regulatória.} A pergunta
  que uma administração enfrenta é se o valor que o modelo devolve corresponde
  ao registro oficial. Construímos cinco tipos de tarefa a partir do trabalho
  cotidiano de um órgão ambiental (Tabela~\ref{tab:tarefas}), cada um atrelado a
  uma fonte oficial, de modo que a correção é adjudicável em vez de uma questão
  de impressão de especialista, e o protocolo se transfere a qualquer domínio
  regulado em que a resposta correta é publicada.
\item \emph{As três correções de localização a que os órgãos recorrem,
  testadas em vez de presumidas.} Enquadramento por persona, prompt em idioma
  nativo e modelo com ajuste regional são, cada um, uma decisão de compra ou de
  fluxo de trabalho. Nenhum deles funciona, e falham de maneiras diferentes:
  perguntar no idioma do próprio país reduz a acurácia em $\nnativa$ pontos
  percentuais, de forma mais acentuada para o idioma de corpus mais escasso; o
  enquadramento de papel equivale a nenhum enquadramento em até dois pontos
  percentuais; e o modelo com ajuste regional recorda o valor vinculante com a
  mesma frequência do modelo-base de que partiu.
\item \emph{Um desenho que separa quem é mal servido de por quê, e relata onde
  não consegue.} Os países são estratificados pela divisão desenvolvido/em
  desenvolvimento da UNCTAD, pela taxonomia de recursos linguísticos de Joshi e
  pelo desenvolvimento, de modo que explicações geográficas, linguísticas e
  econômicas possam ser distinguidas. Entre países, cobertura e desenvolvimento
  são colineares; dentro dos países, o que prevê o déficit de especificidade no
  nível do país é o desenvolvimento, e não a cobertura enciclopédica que o
  protocolo pré-especificou, e os dois não se separam nesta amostra. Documentamos também
  que o proxy de corpus pré-especificado atribui a Wikipédia inglesa a todo
  país da amostra com inglês como língua oficial, de modo que ele acompanha a
  história colonial e não o tamanho do corpus.
\item \emph{Gabarito adjudicado por código, e o que isso muda.} Reconstruir
  duas chaves de tarefa a partir de registros oficiais e mover a comparação de
  faixa para código reduziu o gradiente de desenvolvimento de $\rho=0,51$ para
  $\rho=\nrho$ e mais que dobrou a penalidade do idioma nativo nas mesmas
  respostas; mover a terceira tarefa com valor de registro de um juiz para
  código removeu um confundimento entre juiz e tier. Nesta literatura, a chave
  de gabarito, e não o tamanho da amostra, é a restrição que vincula.
\item \emph{Recursos abertos.} Protocolo, prompts, registros de gabarito com
  suas fontes oficiais, código de análise e o registro de desvios são liberados
  sob CC-BY-4.0 e MIT, com pipeline de reprodução em um comando.
\end{enumerate}
